% TOP_PROBLEM: {"id": 1100902, "title": "Questions in Geometric Group Theory — Q 9.2", "category_id": 17, "category": {"id": 17, "name": "group_theory", "display_name": "Group Theory", "description": "Problems about groups, group actions, representations, and related algebraic structures.", "slug": "group-theory", "order_index": 17, "created_at": "2026-07-31T15:26:25.671Z"}, "status": "open", "problem_number": "AMR-010-0902", "set_id": 13}
% FIRST_POSTED: 2026-10-02
% ENTRY_KIND: statement_only
% UnsolvedMath ID 1100902; source catalogue code AMR-010-0902
% !TeX program = lualatex
% Definition and sources only; this file is not an LLM solution attempt.
\documentclass[11pt,a4paper]{article}
\usepackage[margin=25mm]{geometry}
\usepackage{fontspec}
\setmainfont{lmroman10-regular.otf}[BoldFont=lmroman10-bold.otf,
 ItalicFont=lmroman10-italic.otf,BoldItalicFont=lmroman10-bolditalic.otf]
\usepackage{amsmath,amssymb,mathtools}
\usepackage{unicode-math}
\setmathfont{latinmodern-math.otf}
\AtBeginDocument{\let\setminus\smallsetminus}
\usepackage{xurl,hyperref}
\hypersetup{colorlinks=true,urlcolor=blue}
\setcounter{secnumdepth}{0}
\setlength{\parindent}{0pt}
\setlength{\parskip}{5pt}
\setlength{\emergencystretch}{3em}
\begin{document}
\section*{Questions in Geometric Group Theory — Q 9.2}\label{1100902}
{\small UnsolvedMath ID 1100902 · AMR-010-0902 · Group Theory · AMR Open Problem Lists}
\subsection{Definitions and mathematical statement}
An elementary collapse of a CW complex removes an open cell together with a free codimension-one face; its inverse is an elementary expansion. A simple homotopy equivalence between finite CW complexes is a homotopy equivalence obtained, up to homotopy, from finitely many such moves, with intermediate complexes of unrestricted dimension.

Let $K$ and $L$ be finite two-dimensional CW complexes which are simple homotopy equivalent. Can $K$ be transformed into $L$ by a finite sequence using only expansions and collapses of cell pairs of dimensions $(0,1)$ or $(1,2)$, together with homotopies of the attaching maps of two-cells?

For the last operation, retain the existing one-skeleton and replace an attaching loop
\[
 a:S^1\longrightarrow K^{(1)}
\]
by a loop homotopic to it in the subcomplex consisting of the one-skeleton and the other two-cells. This allows a two-cell to slide over another two-cell while its own attaching map is changed. All intermediate complexes remain at most two-dimensional; subdivision and cellular identifications are harmless changes of CW description.

The problem is the topological Andrews--Curtis question in the source. It does not permit an unrecorded three-cell expansion, since eliminating such higher-dimensional moves is exactly the issue. Nor is ordinary homotopy equivalence alone the hypothesis: the given simple homotopy equivalence excludes a Whitehead-torsion obstruction at the outset.
\subsection{Short English statement}
Can every simple homotopy equivalence between finite two-complexes be realized using only two-dimensional moves and two-cell slides?
\subsection{Sources}
\begin{itemize}
\item[S1] ULAM.AI and the UnsolvedMath Contributors, supplied problems.json, record 1100902 (AMR-010-0902), AMR Open Problem Lists. Catalogue identity and source transcription; CC BY 4.0.
\url{https://huggingface.co/datasets/ulamai/UnsolvedMath}.
\item[S2] Bestvina - Questions in Geometric Group Theory (pdf) (2004), Question 9.2 (PDF page 18). Original source indexed by AMR-010-0902.
\url{https://www.math.utah.edu/~bestvina/eprints/questions-updated.pdf}.
\end{itemize}
\end{document}
